% !TEX root = ../report_template.tex

\section{Discussion}
\label{sec:discussion}
\label{sec:method}
\label{sec:challenges}

The previous section shows that our prototype can compete with
classical baselines on synthetic and real tables within a limited
training budget. However, its performance varies across datasets
and metrics. These results raise two questions: why does training
on synthetic data transfer to real tasks, and where does this
approach still fall short?

TabICLv2 learns to predict unknown labels from labeled examples
across many synthetic tasks. From the prior-data fitted network
(PFN) perspective, it approximates the posterior predictive
distribution induced by the prior: a probability distribution
over unknown labels given the observed examples
\citep{muller2022pfn,qu2026tabiclv2}.
The prior determines the relationships between features and labels
encountered during training. The architecture builds column
representations, combines them into row representations, and uses
in-context learning to predict test labels. On a new dataset, the
labeled rows provide context, while the pretrained model weights
remain fixed.

This approach is interesting because synthetic tasks can be
generated continuously without collecting real labels for every
pretraining task. The same pretrained model can then be applied
to different datasets. However, generating more data does not
necessarily make it representative of real problems. Our dataset
analysis found that 21 of the 57 OpenML datasets have fewer training
rows than the range used in pretraining. Some also contain missing
values or strong class imbalance. Such differences may affect
transfer. The challenge is therefore both to learn from synthetic
tasks and to ensure that these tasks cover important properties
of real data.

Our five-hour run shows that this approach is feasible at the scale
of our project. The model achieved mean ROC AUC values of 0.890 on
prior tables and 0.901 on OpenML, slightly above the strongest
baseline for each source. However, the differences are small and
the confidence intervals overlap, so they do not establish a clear
advantage. Lower log loss and calibration error suggest that the
model produces useful probability estimates within this evaluation.

These strengths do not mean that the model handles every class or
dataset well. On volcanoes-a1, it achieved a high AUC but often
predicted the majority class, resulting in low balanced accuracy.
On hill-valley, logistic regression performed much better than our
prototype. These examples show why several metrics and individual
datasets need to be examined alongside the averages. A mismatch
between the prior and a particular task is one possible
explanation, but our experiments do not isolate the causes of
these failures.

There are also limitations in the evaluation. The best checkpoint
was selected using an OpenML validation split, so results on these
datasets may be optimistic. A second split provides additional
evidence, but it uses the same datasets and is not a fully
independent test. We also evaluated only one training run, and the
baselines were not extensively tuned. The absence of a sustained
rise in validation loss is encouraging, but it does not rule out
overfitting to the validation tasks through checkpoint selection.

With the available compute budget, our next priority would be a
more independent evaluation, followed by targeted experiments.
We would test on datasets not used for checkpoint selection and,
if possible, repeat training with another random seed. We could
then change one aspect of the prior at a time, such as including
smaller contexts, stronger class imbalance, or missing values.
This would help identify which changes actually improve transfer.
Continued pretraining on real data is another possible direction:
\citet{qu2026tabiclv2} discuss evidence that real-data pretraining
can be competitive and can complement synthetic pretraining.
However, this has not been tested in our project. For now, smaller
controlled experiments would help us understand the remaining
weaknesses before committing to a longer training run.
